\documentclass[11pt]{article}
\usepackage[letterpaper,margin=.68in,headheight=16pt,headsep=.18in,footskip=.28in]{geometry}
\usepackage{amsmath,array,enumitem,fancyhdr,hyperref}
\pagestyle{fancy}\fancyhf{}\renewcommand{\headrulewidth}{.3pt}
\fancyhead[L]{\small Bellingham Math Circle / Week 14 / Return visits}
\fancyhead[R]{\small Adult guide}
\fancyfoot[L]{\small RV-14-FAC-v1 / Draft and unpiloted}\fancyfoot[R]{\small\thepage}
\setlength{\emergencystretch}{2em}
\setlength{\parindent}{0pt}\setlength{\parskip}{7pt}
\setlist[itemize]{leftmargin=1.3em,itemsep=3pt,topsep=4pt}
\hypersetup{hidelinks,pdftitle={Week 14 return-visit facilitator guide}}
\begin{document}
{\Large\bfseries Mathematical overview}\par\medskip
\textbf{Facts and limits.} For a triangulated convex n-gon using only original corners and noncrossing diagonals, the graph connecting triangles across shared diagonals is a tree: it is connected, with n-2 triangle nodes and n-3 edges. For n\(\geq\)4 it has at least two leaves. Each leaf is an ear, a triangle with two boundary sides. There may be more than two ears.

The corners admit a three-coloring with all three labels on every triangle, unique up to permuting labels. Remove ears to a single triangle, color it, then insert ears in reverse order: the new ear corner gets the third label opposite its two already colored neighbors. The smallest color class touches every triangle and contains at most \(\lfloor n/3\rfloor\) corners. This is a combinatorial cover, not physical visibility in a room.

Of the 14 triangulations of a regular labeled hexagon, six have half-turn symmetry, two have one-third-turn symmetry and none have one-sixth-turn symmetry. Labels remain fixed when counting different drawings; symmetry tests whether a rotation carries a diagonal set back to itself.

\textbf{Children's destination.} Peel a structural object into smaller pieces; use a coloring to certify a triangle cover; find a complete symmetry collection by orbit constraints rather than guesswork.

\textbf{Entry map.} Problem 1 (student pp. 1--2): K--1 and up, recognize two outside sides and erase a corner with adult demonstration. Problem 2 (student pp. 3--4): grades 2--3 and up, label corners and check every triangle. Problem 3 (student pp. 5--6): grades 4--5, turning/tracing, fractions of a turn and accounting for all possibilities. Formal tree or group language is optional.

\textbf{New versus base.} The base counts triangulations and studies flips and fans. These pages add ears/adjacency, corner coloring/coverage, and geometric symmetry.
\newpage {\Large\bfseries Prepare a return visit}\par\medskip
\textbf{Status and use.} Three investigations extend one familiar theme. These companions are separate from the base packets and may support several visits. All pages are new and unpiloted. Printed labels describe prerequisites, not age restrictions. Offer one investigation’s pages at a time; completing every page is not the goal.

\textbf{Materials and preparation.} For each pair, pencil/eraser, ruler, three R/B/Y pencils or letter labels and two tracing sheets. Reserve six sets for eleven children and at least twelve spare sheets. Main polygons are about four inches across, with fixed corner labels and equal x/y scaling. No tile fit is required. A turn must keep the outline fixed; tracing lets children check the geometric symmetry while preserving the corner names.

\textbf{Staffing.} Current context: eleven children, fixed KK11 / 3333 / 445 tables, one adult anchored to each table. Use pairs; the upper third child rotates as reader or checker. Routine adult reading, arrow drawing or recording can leave the mathematical decisions with children. If adults must choose every move, use the earlier entry rather than run the investigation for them.

\textbf{Short common launch.} Let children draw a triangle division on spare paper. Point to a triangle with two outline sides and cover it with a small paper flap; then show how erasing its middle corner and the two short outline sides leaves a smaller polygon bounded by the third side. Demonstrate a turn on a square diagonal as a non-task example. At the rotation table, one, two and three corner steps on an empty regular hexagon show one-sixth, one-third and one-half turns if needed. Do not announce the minimum number of ears or the hexagon symmetry counts.

\textbf{Suggested first route.} Supply the two-page pair for the chosen investigation together: pp. 1--2, 3--4 or 5--6. Start with finding and removing ears. Offer corner coloring after children can reliably identify every triangle. Return for rotation symmetry when a child can compare two diagonal sets with fixed corners; tracing or adult label reading may help without making the choices.

\textbf{Flexible timing.} Allow 3--5 minutes of material play and 2--4 minutes for the common demonstration, then 15--25 minutes for one investigation and 5 minutes to share. A longer second visit can spend 30--45 minutes on the same investigation, including unfinished examples or its explanation. These are planning estimates, not observed timings. Do not require all three within one hour.

\textbf{Questions and hints.} Start with ``What can you try?'' and ``What would count as success?'' Check the actual rule before giving a strategy. Optional questions and hints are keyed below; use them only after a real attempt. A counter argument or drawing may be a complete explanation.

\textbf{Source and pilot record.} Mathematical examples and arguments here are original local follow-ups to the current base theme. Consulted teaching source: \emph{Math Circle by the Bay}, Preface pp. vii--x (local PDF pp. 8--11), on interaction, manipulatives, hints, explanations and themes spanning multiple years. Our exact launch, entry route and timing are design inferences. Year-1 \emph{Handouts 6}, Problems 6.2--6.6, provided a local example of connecting representations. Record actual problem, starting route, examples tried, what needed adult rescue, and what children want to resume. Distinguish observations from hypotheses. Counter fit, materials stock, procedures and classroom response have not been physically tested.
\newpage
\textbf{Problem 1: ears and the triangle tree.} An ear has two polygon boundary sides. In a hexagon triangulated by the alternating central triangle, the three outside triangles are all ears. A fan has exactly two ears. Thus ``at least two'' is sharp, but ``exactly two'' is false.

The printed task asks whether different choices leave different last triangles. Yes: on the fan, peel B,C,D to leave AEF, or F,E,D to leave ABC. On the central drawing, peel B,D,F to leave ACE, or D,F,E to leave ABC. Every original triangle can be the last one: retain its node in the adjacency tree and repeatedly remove a leaf away from it. The fan has eight complete corner-removal words, and the central drawing twelve; these are adult checks, not required child catalogs.

Put one point in every triangle and connect points across every shared inside diagonal. This graph is connected: adjacent cells can be reached by crossing successive diagonals inside the polygon. It has n-2 nodes and n-3 edges, so it is a tree. A finite tree with at least two nodes has at least two leaves; otherwise a longest simple path could be extended at an endpoint. A leaf triangle has only one neighboring triangle and therefore two boundary sides: it is an ear.

\textbf{A direct peeling entry.} Remove an ear by erasing its middle polygon corner. Its third side becomes boundary, leaving a triangulated convex (n-1)-gon. Continue until one triangle remains. Children can record just the corner names erased; the original drawing lets them rebuild in reverse.

\textbf{If needed.} Trace the outline sides of one candidate triangle. A triangle with only one outside side is not an ear, even if it is narrow or near an edge.

\textbf{Extension.} What is the greatest number of ears an n-gon can have? No two ear middle corners can be adjacent when n\(\geq\)5, so at most \(\lfloor n/2\rfloor\); alternating ear constructions attain it. Treat the quadrilateral separately: its two ears have nonadjacent middle corners. A new drawing and bound are a worthwhile return task.
\newpage
\textbf{Problem 2: a corner cover from three labels.} Color a final triangle R,B,Y. Reinsert ears in reverse deletion order. Each added corner is adjacent to the two ends of its ear diagonal, which already have different labels; the third label is forced. This proves existence and uniqueness up to the initial label permutation without assuming an arbitrary traversal always encounters an uncolored corner.

Every triangle has all three labels, so choosing every R corner, or every B corner, or every Y corner touches every triangle. The smallest class has at most \(\lfloor n/3\rfloor\) members, because the three class sizes sum to n.

For the hexagon with central triangle ACE and ears ABC,CDE,EFA, one possible cyclic labeling is R,Y,B,R,Y,B. Each label has two corners, and two corners suffice to touch every triangle. One corner cannot suffice: no corner belongs to all three outside ears. The six minimum pairs are \(AC,AD,AE,BE,CE,CF\); mixed-color selections are valid too. For a fan at A, the unique minimum is A alone, and its color class has just that one corner; it touches every triangle. Coloring supplies a reliable upper bound; an arbitrary chosen class need not be a minimum cover.

\textbf{If needed.} Cover all but one already labeled triangle. Ask which label its remaining corner needs. Let children choose their covering corners and test each triangle before offering the smallest-color-class idea.

\textbf{Extension.} Build a triangulation requiring more than one chosen corner, then seek a cover smaller than the most obvious color class. Keep ``touch a triangle at a vertex'' explicit; covering the drawing with a big physical counter is not the rule.
\newpage
\textbf{Problem 3: turn-symmetric fillings.} Using the printed corner names A–F cyclically, the six half-turn-symmetric diagonal sets are
\[
\begin{array}{ll}
AC,AD,DF&AC,CF,DF\\
AD,BD,AE&BD,BE,AE\\
BE,BF,CE&BF,CE,CF.
\end{array}
\]
(These use compact adult endpoint notation; check drawings rather than teach it as an extra convention.) The two one-third-turn sets are \(AC,CE,AE\) and \(BD,DF,BF\). No one-sixth-turn filling exists.

\textbf{Why the half-turn list is complete.} A triangulation has three diagonals. Under a half-turn, nonfixed diagonals occur in pairs, so at least one diagonal is fixed. Such a diagonal must be a diameter. Two diameters would cross, so there is exactly one. Choose its three possible positions, then triangulate one of its two quadrilaterals in either of two ways. The other half is forced: \(3\cdot2=6\).

\textbf{One-third and one-sixth turns.} A one-third-turn orbit of a diameter consists of three crossing diameters and is forbidden. The other possible diagonal orbits are the two alternating central triangles, both legal. A one-sixth turn would require either all three diameters (crossing) or an orbit of six short diagonals (more than the available three); neither works.

\textbf{If needed.} Give a tracing copy. Turn the diagonals, keep the outline aligned and compare each line, rather than rotating the labels and deciding that every picture is automatically ``the same.''

\textbf{Checks.} The independent source audit enumerates noncrossing three-diagonal sets, all triangle cells, ears, valid colorings, corner covers and rotation-fixed sets. It confirms 14 total, six half-turn, two one-third, zero one-sixth.
\end{document}
